
homomorphisme se factorise en :
@@M0@@
et le premier homomorphisme est injectif. Comme, d'après VI 4.3, 4.6 et 4.7, @@M1@@ est libre sur @@M2@@, cela entraîne l'injectivité de :
@@M3@@

Soit donc @@M4@@. On procède par récurrence sur le rang de @@M5@@. Le cas @@M6@@ étant trivial, supposons la propriété prouvée pour @@M7@@. On introduit alors @@M8@@, et le noyau @@M9@@ de l'homomorphisme canonique @@M10@@, où @@M11@@ est le morphisme structural de @@M12@@. Alors @@M13@@ s'identifie au schéma des drapeaux de @@M14@@ sur @@M15@@, et on voit d'après l'hypothèse de récurrence qu'il suffit de prouver 3.1 pour @@M16@@. Compte tenu de VI 1.1, le résultat à démontrer peut donc s'énoncer :

\textbf{Corollaire 3.2.} \uline{Soient} @@M17@@ \uline{un schéma noethérien}, @@M18@@ \uline{un} @@M19@@\uline{-Module localement libre de rang} @@M20@@, @@M21@@ \uline{le fibré projectif associé}, @@M22@@ \uline{le morphisme structural de} @@M23@@, \uline{et} @@M24@@ \uline{la classe dans} @@M25@@ \uline{du faisceau} @@M26@@. \uline{Alors l'homomorphisme} :
@@M27@@
\uline{défini par} :
@@M28@@
\uline{est un isomorphisme}.

On a déjà vu que cet homomorphisme est surjectif, puisque les @@M29@@, @@M30@@, forment une base de @@M31@@ sur @@M32@@; pour voir qu'il est injectif, on utilise l'homomorphisme @@M33@@, défini car @@M34@@ est projectif.
